\begin{definition}[Условие Татта]
\(\odd(G)\) — число нечётных компонент связности графа \(G\).
Граф удовлетворяет \emph{условию Татта}, если \(\odd(G-S)\le|S|\) для всех
\(S\subseteq V\). Множество \(S\) с \(\odd(G-S)>|S|\) называется
\emph{множеством Татта}.
\end{definition}
\begin{theorem}[Татт, 1947]
В графе \(G\) есть совершенное паросочетание тогда и только тогда, когда
\(\odd(G-S)\le|S|\) для всех \(S\subseteq V\).
\end{theorem}
\textit{Необходимость.} Пусть \(M\) — совершенное паросочетание и \(S\subseteq V\).
Нечётную компоненту \(Q\) графа \(G-S\) нельзя разбить на пары внутри себя,
значит, некоторая её вершина сопоставлена ребром \(M\) вершине вне \(Q\);
другие компоненты \(G-S\) с \(Q\) не смежны, поэтому партнёр лежит в \(S\).
Партнёры для разных компонент различны, так что \(\odd(G-S)\le|S|\).
\begin{graphillustration}{Множество Татта: после удаления $S=\{s\}$ остаются три нечётные компоненты, а партнёров в $S$ хватит только одной. Совершенного паросочетания нет.}{fig:04-tutte-set}
\begin{tikzpicture}[font=\small]
\node[gi selected] (s) at (0,0) {$s$};
\node[gi vertex] (a) at (-2.2,-1.4) {};
\node[gi vertex] (b1) at (0,-1.4) {};
\node[gi vertex] (b2) at (-0.6,-2.4) {};
\node[gi vertex] (b3) at (0.6,-2.4) {};
\node[gi vertex] (c) at (2.2,-1.4) {};
\draw[gi edge] (s) -- (a); \draw[gi edge] (s) -- (b1); \draw[gi edge] (s) -- (c);
\draw[gi edge] (b1) -- (b2) -- (b3) -- (b1);
\node[gi group,fit=(a)] {};
\node[gi group,fit=(b1)(b2)(b3)] {};
\node[gi group,fit=(c)] {};
\node[gi note,anchor=west] at (3.2,-1.2) {$\odd(G-S)=3>1=|S|$};
\end{tikzpicture}
\end{graphillustration}


\textit{Достаточность} (доказательство Ловаса). Две вспомогательные леммы.
\begin{lemma}\label{lem:odd-edge}
Если \(uv\notin E\), то \(\odd(G+uv)\le\odd(G)\).
\end{lemma}
\begin{proof}
Если \(u,v\) в одной компоненте, компоненты не меняются. Иначе две компоненты
сливаются: чёт\(+\)чёт \(=\) чёт, нечёт\(+\)чёт \(=\) нечёт (число нечётных
не меняется), нечёт\(+\)нечёт \(=\) чёт (уменьшается на 2).
\end{proof}
\begin{lemma}
Если \(G\) удовлетворяет условию Татта, то и \(G+uv\) удовлетворяет.
\end{lemma}
\begin{proof}
Для \(S\subseteq V\): если \(u\in S\) или \(v\in S\), то \((G+uv)-S=G-S\);
иначе \((G+uv)-S=(G-S)+uv\), и по лемме~\ref{lem:odd-edge}
\(\odd\bigl((G+uv)-S\bigr)\le\odd(G-S)\le|S|\).
\end{proof}
\begin{proof}[Доказательство достаточности]
Предположим, что теорема неверна. Среди контрпримеров (условие Татта
выполнено, совершенного паросочетания нет) на данном множестве вершин \(V\)
выберем граф \(G\) с \emph{максимальным} числом рёбер. Тогда
\begin{itemize}\tightlist
\item \(n\) чётно: при \(S=\varnothing\) имеем \(\odd(G)\le0\);
\item для любого \(e\notin E\) граф \(G+e\) удовлетворяет условию Татта и имеет
  больше рёбер, значит, \emph{в \(G+e\) есть совершенное паросочетание};
\item \(G\ne K_n\) (в \(K_n\) с чётным \(n\) совершенное паросочетание есть).
\end{itemize}
Пусть \(U=\{u\in V: d(u)=n-1\}\) — вершины, смежные со всеми; \(U\ne V\).

\begin{lemma}\label{lem:cliques}
\(G-U\) — дизъюнктное объединение клик.
\end{lemma}
\textit{Доказательство леммы.} Иначе в одной компоненте \(G-U\) найдутся \(x,y,z\) с \(xy,yz\in E\), \(xz\notin E\)
(возьмём кратчайший путь между несмежными вершинами компоненты). Так как
\(y\notin U\), есть \(w\ne y\) с \(yw\notin E\). По выбору \(G\) существуют совершенные
паросочетания \(M_1\) графа \(G+xz\) и \(M_2\) графа \(G+yw\); поскольку в \(G\)
совершенного паросочетания нет, \(xz\in M_1\) и \(yw\in M_2\).
Граф \(H=(V,M_1\mathbin{\triangle}M_2)\) состоит из изолированных вершин и
чётных чередующихся циклов (лемма билета~\ref{ticket:1}; путей нет: оба
паросочетания совершенны), причём \(xz\) и \(yw\) лежат на циклах.

\emph{Случай 1: \(xz\) и \(yw\) на разных циклах \(C_1\), \(C_2\).} Возьмём рёбра
\(M_2\) на \(C_1\), рёбра \(M_1\) на \(C_2\) и рёбра \(M_1\) на остальных вершинах.
Это совершенное паросочетание без \(xz\) и \(yw\), т.\,е. в \(G\) — противоречие.

\emph{Случай 2: на одном цикле \(C\).} Пойдём по \(C\) от \(y\) через ребро \(yw\);
не умаляя общности (\(x\) и \(z\) симметричны: обе смежны с \(y\)), раньше
встретим \(z\), так что \(C=y\,w\ldots z\,x\ldots y\). Путь
\[
P=\underbrace{x\ldots y}_{\text{дуга }C}\;\; yz\;\;\underbrace{z\ldots w}_{\text{дуга }C}
\]
проходит все вершины \(C\) и не использует \(xz\), \(yw\); его рёбра лежат в \(G\).
Число вершин \(C\) чётно, поэтому рёбра \(P\) через одно дают совершенное
паросочетание \(V(C)\). Добавив \(M_1\) вне \(C\), получим совершенное паросочетание
\(G\) — противоречие. Лемма доказана.
\begin{graphillustration}{Случай 2 леммы о кликах. Слева цикл $C$ из $M_1\triangle M_2$: толстые рёбра — $M_1$, штриховые — $M_2$; $xz\in M_1$ и $yw\in M_2$ — добавленные рёбра. Справа путь $P$: дуга $x\ldots y$, ребро $yz\in E$, дуга $z\ldots w$; рёбра $P$ через одно (толстые) дают совершенное паросочетание на $V(C)$.}{fig:04-cycle}
\begin{tikzpicture}[font=\small]
\begin{scope}
\foreach \i/\a/\l in {0/90/y,1/30/w,2/-30/p,3/-90/z,4/-150/x,5/150/q}
  \node[gi vertex] (c\i) at (\a:1.4) {$\l$};
\draw[gi alternate] (c0) -- node[above right,inner sep=1pt] {\scriptsize новое} (c1);
\draw[gi match] (c1) -- (c2);
\draw[gi alternate] (c2) -- (c3);
\draw[gi match,dash pattern=on 4pt off 2pt] (c3) -- node[below,inner sep=2pt] {\scriptsize новое} (c4);
\draw[gi alternate] (c4) -- (c5);
\draw[gi match] (c5) -- (c0);
\draw[gi edge,densely dotted] (c0) -- (c3);
\node[gi note] at (0,-2.2) {$C$; пунктир $yz\in E(G)$};
\end{scope}
\node[gi note] at (3.0,0) {$\Longrightarrow$};
\begin{scope}[shift={(6,0)}]
\foreach \i/\a/\l in {0/90/y,1/30/w,2/-30/p,3/-90/z,4/-150/x,5/150/q}
  \node[gi vertex] (d\i) at (\a:1.4) {$\l$};
\draw[gi match] (d4) -- (d5);
\draw[gi edge] (d5) -- (d0);
\draw[gi match] (d0) -- (d3);
\draw[gi edge] (d3) -- (d2);
\draw[gi match] (d2) -- (d1);
\draw[gi removed] (d0) -- (d1);
\draw[gi removed] (d3) -- (d4);
\node[gi note] at (0,-2.2) {$P=x\,q\,y\,z\,p\,w$};
\end{scope}
\end{tikzpicture}
\end{graphillustration}


\emph{Завершение.} По условию \(\odd(G-U)\le|U|\). В каждой чётной клике \(G-U\)
разобьём вершины на пары. В каждой нечётной клике разобьём на пары все вершины,
кроме одной, а её сопоставим своей вершине \(U\) (вершины \(U\) смежны со всеми, а
различных вершин хватает). Осталось чётное число вершин \(U\) (всего \(n\) чётно,
а покрыто уже чётное число); они попарно смежны, разобьём их на пары.
Получено совершенное паросочетание \(G\) — противоречие.
\end{proof}
\begin{graphillustration}{Завершение доказательства Татта: $G-U$ — несвязные клики. Чётные клики разбиваются на пары внутри; в каждой нечётной одна вершина уходит к своей вершине $U$; остаток $U$ — клика с чётным числом вершин.}{fig:04-finish}
\begin{tikzpicture}[font=\small]
\node[gi vertex] (u1) at (0,2) {};
\node[gi vertex] (u2) at (1.2,2) {};
\node[gi vertex] (u3) at (2.4,2) {};
\node[gi vertex] (u4) at (3.6,2) {};
\node[gi group,fit=(u1)(u2)(u3)(u4),label=above:{$U$ (смежны со всеми)}] {};
\draw[gi match] (u3) -- (u4);
\node[gi vertex] (a1) at (-0.6,0) {};
\node[gi vertex] (a2) at (0.4,0) {};
\node[gi vertex] (a3) at (-0.1,-0.8) {};
\draw[gi edge] (a1) -- (a2) -- (a3) -- (a1);
\draw[gi match] (a2) -- (a3);
\node[gi vertex] (b1) at (1.8,0) {};
\draw[gi match] (a1) -- (u1);
\draw[gi match] (b1) -- (u2);
\node[gi vertex] (c1) at (3.0,0) {};
\node[gi vertex] (c2) at (4.0,0) {};
\node[gi vertex] (c3) at (3.0,-0.9) {};
\node[gi vertex] (c4) at (4.0,-0.9) {};
\draw[gi edge] (c1) -- (c2) -- (c4) -- (c3) -- (c1) -- (c4); \draw[gi edge] (c2) -- (c3);
\draw[gi match] (c1) -- (c2); \draw[gi match] (c3) -- (c4);
\node[gi group,fit=(a1)(a2)(a3)] {};
\node[gi group,fit=(b1)] {};
\node[gi group,fit=(c1)(c2)(c3)(c4)] {};
\node[gi note,anchor=west] at (4.8,-0.4) {клики $G-U$:\\две нечётные, одна чётная};
\end{tikzpicture}
\end{graphillustration}


\begin{cor}[Петерсен]
Любой 3-регулярный граф без мостов имеет совершенное паросочетание.
\end{cor}
\begin{proof}[Набросок (по книге)]
Пусть \(S\subseteq V\), \(Q\) — нечётная компонента \(G-S\). Сумма степеней в \(Q\)
равна \(3|Q|\) — нечётна, а внутренние рёбра дают чётный вклад, поэтому
из \(Q\) в \(S\) ведёт нечётное число рёбер, т.\,е. не меньше \(3\) (одно ребро было бы
мостом). Итого \(3\odd(G-S)\le\) рёбер, входящих в \(S\), \(\le3|S|\).
\end{proof}
